\documentclass[
10pt, % Main document font size
a4paper, % Paper type, use 'letterpaper' for US Letter paper
oneside, % One page layout (no page indentation)
%twoside, % Two page layout (page indentation for binding and different headers)
headinclude,footinclude, % Extra spacing for the header and footer
BCOR5mm, % Binding correction
]{scrartcl}
%%%%%%%%%%%%%%%%%%%%%%%%%%%%%%%%%%%%%%%%%
% Arsclassica Article
% Structure Specification File
%
% This file has been downloaded from:
% http://www.LaTeXTemplates.com
%
% Original author:
% Lorenzo Pantieri (http://www.lorenzopantieri.net) with extensive modifications by:
% Vel (vel@latextemplates.com)
%
% License:
% CC BY-NC-SA 3.0 (http://creativecommons.org/licenses/by-nc-sa/3.0/)
%
%%%%%%%%%%%%%%%%%%%%%%%%%%%%%%%%%%%%%%%%%

%----------------------------------------------------------------------------------------
%	REQUIRED PACKAGES
%----------------------------------------------------------------------------------------

\usepackage[
nochapters, % Turn off chapters since this is an article        
beramono, % Use the Bera Mono font for monospaced text (\texttt)
eulermath,% Use the Euler font for mathematics
pdfspacing, % Makes use of pdftex’ letter spacing capabilities via the microtype package
dottedtoc % Dotted lines leading to the page numbers in the table of contents
]{classicthesis} % The layout is based on the Classic Thesis style

\usepackage{arsclassica} % Modifies the Classic Thesis package
\usepackage{ulem}
\usepackage{graphicx}
\usepackage{float}
\usepackage{algorithm2e}
\usepackage[export]{adjustbox}
\usepackage[utf8]{inputenc}  
\usepackage[english,frenchb]{babel} 
\usepackage[T1]{fontenc} % Required for accented characters
\usepackage{multicol}
\usepackage{graphicx} % Required for including images
\graphicspath{{Figures/}} % Set the default folder for images

\usepackage{enumitem} % Required for manipulating the whitespace between and within lists

\usepackage{lipsum} % Used for inserting dummy 'Lorem ipsum' text into the template


\expandafter\def\csname ver@subfig.sty\endcsname{}


\usepackage{svg}
\usepackage{hyperref}

\usepackage{amsmath,amssymb,amsthm} % For including math equations, theorems, symbols, etc

\usepackage{varioref} % More descriptive referencing

%----------------------------------------------------------------------------------------
%	THEOREM STYLES
%---------------------------------------------------------------------------------------

\theoremstyle{definition} % Define theorem styles here based on the definition style (used for definitions and examples)
\newtheorem{definition}{Definition}

\theoremstyle{plain} % Define theorem styles here based on the plain style (used for theorems, lemmas, propositions)
\newtheorem{theorem}{Theorème}
\newtheorem{lemma}[theorem]{Lemme}
\newtheorem{proposition}[theorem]{Proposition}

\theoremstyle{remark} % Define theorem styles here based on the remark style (used for remarks and notes)

%----------------------------------------------------------------------------------------
%	HYPERLINKS
%---------------------------------------------------------------------------------------

\hypersetup{
%draft, % Uncomment to remove all links (useful for printing in black and white)
colorlinks=true, breaklinks=true, bookmarks=true,bookmarksnumbered,
urlcolor=webbrown, linkcolor=RoyalBlue, citecolor=webgreen, % Link colors
pdftitle={}, % PDF title
pdfauthor={\textcopyright}, % PDF Author
pdfsubject={}, % PDF Subject
pdfkeywords={}, % PDF Keywords
pdfcreator={pdfLaTeX}, % PDF Creator
pdfproducer={LaTeX with hyperref and ClassicThesis} % PDF producer
}

\hyphenation{Fortran hy-phen-ation} 

\title{\normalfont{FTML Exercices 1}\\
\large{Pour le 23 avril 2026}
} 
\author{\spacedlowsmallcaps{}
} 
\date{}
\begin{document}
\renewcommand{\sectionmark}[1]{\markright{\spacedlowsmallcaps{#1}}}
\lehead{\mbox{\llap{\small\thepage\kern1em\color{halfgray} \vline}\color{halfgray}\hspace{0.5em}\rightmark\hfil}} 
\pagestyle{scrheadings}
\maketitle 
\setcounter{tocdepth}{3} 

\tableofcontents


\section{\large\color{Blue}Probabilités et statistiques}

\subsection{\color{MidnightBlue}P1}

\subsubsection{Enoncé}
\label{par:enoncé}

Calculer l'espérance et la variance des variables aléatoires réelles suivantes.
\begin{itemize}
    \item $X_1$ de loi uniforme sur $[0, 1]$.
    \item $X_2$ de loi uniforme sur $[-1, 1]$.
\end{itemize}


\subsection{\color{MidnightBlue}P2}

\subsubsection{Enoncé}
\label{par:enoncé}

Calculer l'espérance et la matrice de variance-covariance du vecteur aléatoire suivant

\begin{equation}
    Y = (Y_1, Y_2)
\end{equation}

où
\begin{itemize}
    \item $Y_1$ suit une loi de Bernoulli de paramètre $p$
    \item $Y_2$ suit une loi normale $ \mathcal{N}(\mu, \sigma)$.
    \item On suppose $Y_1$ et $Y_2$ indépendantes.
\end{itemize}

Le théorème de transfert est souvent utilisé pour réaliser ce type de calcul.

\url{https://en.wikipedia.org/wiki/Law_of_the_unconscious_statistician} 

Cependant, il ne faut pas oublier qu'il s'agit d'un théorème et non d'une définition.

\subsection{\color{MidnightBlue}P3}

\subsubsection{Enoncé}
\label{par:enoncé}

Calculer l'espérance et la matrice de variance-covariance du vecteur aléatoire suivant

\begin{equation}
    Z = (Z_1, Z_2)
\end{equation}

où
\begin{itemize}
    \item $Z_1$ suit une loi uniforme sur $[1, 2]$
    \item $Z_2=Z_1^2$.
\end{itemize}


\subsection{\color{MidnightBlue}P4}

\subsubsection{Enoncé}
\label{par:enoncé}

Le système musical occidental contient 12 notes. Un mode heptatonique contient 7 notes. Si je joue nombre $n$ de notes au hasard, et de façon indépendante et équiprobable, au bout de combien de notes ai-je joué au moins une note hors du mode avec une probabilité supérieure à 0.9 ?

\section{\large\color{Blue}Expected value as a minimization}

\subsubsection{Enoncé}%
\label{ssub:Enoncé}

Soit $X$ une variable aléatoire réelle ayant un moment d'ordre $2$. Montrer que
son espérance $E(X)$ est la quantité minimisant la fonction de variable réelle $t\mapsto E\big( (X-t)^2 \big)$

\section{\large\color{Blue}Calcul différentiel}

Connaître et être à l'aise avec les définitions des gradients et jacobiennes d'applications. Bien noter qu'un gradient est simplement la transposée d'une jacobienne, et qu'il n'est défini que lorsque la fonction est à valeurs scalaires.

Nous ferons des rappels en cours sur les différentielles et les gradients d'applications, mais voici quelques exercices pour s'habituer à les manipuler. Vous pouvez trouver la définition du gradient à la section 1.3 de lecture\textunderscore notes.pdf ou utiliser d'autres sources qui vous conviennent.
\\

\url{https://fr.wikipedia.org/wiki/Gradient} 

\url{https://fr.wikipedia.org/wiki/Diff%C3%A9rentielle} 
\\

\subsection{\color{MidnightBlue}C1}

\subsubsection{Enoncé}
\label{par:enoncé}


Calculer le gradient en tout point des applications suivantes.

$$
f_1 = \left\{
    \begin{array}{ll}
        \mathbb{R}^2 \rightarrow \mathbb{R}  \\
	(x,y)\mapsto 3
    \end{array}
\right.
$$

$$
f_2 = \left\{
    \begin{array}{ll}
        \mathbb{R}^2 \rightarrow \mathbb{R}  \\
	(x,y)\mapsto x^3+\sin y
    \end{array}
\right.
$$

$$
f_3 = \left\{
    \begin{array}{ll}
        \mathbb{R}^2 \rightarrow \mathbb{R}  \\
	(x,y)\mapsto x^3\sin y
    \end{array}
\right.
$$

$$
f_4 = \left\{
    \begin{array}{ll}
        \mathbb{R}^d \rightarrow \mathbb{R}  \\
	x\mapsto ||x||_2^2
    \end{array}
\right.
$$

où $||.||_2$ est la norme euclidienne dans $\mathbb{R}^d$.

\end{document}
